\section{Conclusion}
In this work, we bring attention to the \textit{counterfeit samples} of a code language model: incorrect programs that a model thinks are correct and can pass surface-level correctness checks. We observe that in a sense, these counterfeit samples are adversarial to the model: models often cannot assess their correctness, reason about their execution, and struggle to repair them. Compared to other models, GPT-4 may be different from other evaluated models in this regard, in that they are much less susceptible to the traps we observe on counterfeit samples from other models.

While we operate in the domain of code, where it is simple to precisely check a model's understanding, we suspect that the same phenomena occur more generally in language models, which is consistent with the findings from \citet{west2023generative}. Because models being able to understand their own counterfeit samples is a prerequisite to strong self-repair and self-verification schemes, we recommend that others be critical and careful in light of our findings.


